\documentclass{article}
\usepackage{amsmath,amssymb}
\begin{document}
\title{Five-band saturation: exact equations and a conditional obstruction}
\author{}\date{}\maketitle

No infinite family saturating the second-order dimension bound is proved
here. These notes identify the remaining nonlinear conditions and rule out
one restricted asymptotic mechanism.

\paragraph{Exact five-band equations.}
Let $N=5k-1$ and consider $k$ vectors with disjoint supports in the Dicke
residue classes modulo $k$. Write $c_r$ for the coefficient in the unique
vector containing excitation $r$, and $w_r=|c_r|^2$. Each residue class
must satisfy
\[
 \sum_{r\equiv a\ (k)} w_r=1,\quad
 \sum_{r\equiv a\ (k)} (N/2-r)w_r=0,\quad
 \sum_{r\equiv a\ (k)} (N/2-r)^2w_r=\frac{N(N+2)}{12}.
\]
For positive coefficients, define the edge variables
\[
 Y_r=\sqrt{(r+1)(N-r)}\,\overline{c_{r+1}}c_r,
 \qquad 0\le r<N.
\]
Their moduli obey
\[
 |Y_r|^2=(r+1)(N-r)w_rw_{r+1}.
\]
The remaining first- and second-rank compression conditions are exactly
\[
 \sum_{\substack{r\equiv a\ (k)\\r<N}}Y_r=0,\qquad
 \sum_{\substack{r\equiv a\ (k)\\r<N}}(N/2-r-1/2)Y_r=0,
\]
and
\[
 \sum_{\substack{r\equiv a\ (k)\\r<N-1}}
 \frac{Y_rY_{r+1}}{w_{r+1}}=0.
\]
All shifts near residue boundaries are included by these global indices.
Phases can be reconstructed successively from the $Y_r$ along the physical
chain of excitations. The last equation explains why independently closing
polygons for neighboring pairs does not settle second-order anticoherence:
it involves products of successive edge phases.

\paragraph{Obstruction to a smooth fixed profile.}
Suppose a proposed asymptotic five-band family has neighboring residue
weights and phases converging to the same profile, so that the continuum
versions of its first- and second-step equations use one azimuth per band.
Its five band weights and azimuths must then form a weighted spherical
2-design at the prescribed latitudes.

Every node of any positive weighted spherical 2-design has weight at most
$1/4$. Indeed, the four rows of the matrix with columns
\[
 \sqrt{w_l}\,(1,\sqrt3\,\mathbf n_l)^T
\]
are orthonormal. The squared norm of a column is $4w_l$, a diagonal entry
of an orthogonal projector, so $4w_l\le1$.

At the endpoint of the five-band residue interval the limiting latitudes
are
\[
 1,\quad\frac35,\quad\frac15,\quad-\frac15,\quad-\frac35.
\]
With weights summing to one and bounded above by $1/4$, the smallest possible
weighted mean is zero, attained uniquely by
\[
 (w_0,w_1,w_2,w_3,w_4)=(0,1/4,1/4,1/4,1/4).
\]
The associated second moment is $1/5$, whereas a spherical 2-design requires
$1/3$. Hence this smooth fixed-profile mechanism cannot yield saturation
throughout the five-band interval.

This is not an exclusion theorem for general AC subspaces or even for all
five-band constructions. Neighboring residues may oscillate rather than
approach a common smooth profile. Such oscillations, or a different support
structure, remain possible.
\end{document}
